\documentclass[11pt]{article}
\usepackage[T1]{fontenc}
\usepackage[margin=1in]{geometry}
\usepackage{amsmath,amssymb,amsthm,booktabs,array}
\usepackage{hyperref}
\hypersetup{colorlinks=true,urlcolor=blue,citecolor=blue,linkcolor=blue,
  pdftitle={Residual order four in GP(7,2): an alternative proof for Hausmann's Question 5.4 (arXiv:2611.05178v1, Section 5), after AFTD (2026)},
  pdfauthor={Haobo Ma and Wenlin Zhang}}
\newtheorem{theorem}{Theorem}
\newtheorem{lemma}{Lemma}
\newtheorem{corollary}{Corollary}
\newcommand{\GP}{\operatorname{GP}}
\newcommand{\cc}{\operatorname{cc}}
\title{Residual order four in $\GP(7,2)$:\\
  an alternative proof for Hausmann's Question 5.4\\
  \large arXiv:2611.05178v1, Section 5; after AFTD (2026)}
\author{Haobo Ma\\\small ChronoAI Pte Ltd, Singapore
  \and Wenlin Zhang\\\small National University of Singapore\\
  \small ChronoAI Pte Ltd, Singapore}
\date{November 9, 2026}
\begin{document}
\maketitle
\begin{abstract}
Following AFTD's negative answer dated November 5, 2026, to Hausmann's
Question 5.4 in \emph{Weighted Max Cut in Subcubic Graphs}
(arXiv:2611.05178v1, Section 5), we give a different counterexample
with an exact residual-order calculation: the generalized
Petersen graph $\GP(7,2)$ has odd girth five, but deleting any independent
set leaves a component with at least four vertices. We prove this by
classifying its maximal independent sets into five symmetry classes.
The bound four is attained on this graph.
The same classification gives an obstruction to the nonempty-independent-side
version of Question 5.3; its printed quantifiers admit an empty-side solution.
\end{abstract}

\section{The question and the counterexample}
Hausmann asks whether every finite simple subcubic graph of odd girth
exactly five admits a partition $V(G)=L\mathbin{\dot\cup}R$ with $L$
independent and every connected component of $G[R]$ of order at most
three \cite[Question 5.4]{hausmann}. Here order counts vertices, including
isolated vertices, and odd girth means the length of a shortest odd cycle.
The graph below answers this question negatively even within connected cubic graphs.
AFTD \cite{aftd} previously gave the eight-vertex Wagner graph as a
counterexample. Our alternative uses $\GP(7,2)$ and determines its exact
residual component-order optimum by a complete five-class argument.

For a nonempty finite graph define
\[
 \kappa(G)=\min_{L\text{ independent}}
       \max_{C\in\cc(G[V(G)\setminus L])}|V(C)|,
\]
where the maximum is zero for an empty complement.
Let $H=\GP(7,2)$ have vertices $a_i,b_i$, indexed modulo seven, and edges
\begin{equation}\label{eq:edges}
 a_i a_{i+1},\qquad a_i b_i,\qquad b_i b_{i+2}
 \quad(i\in\mathbb Z/7\mathbb Z).
\end{equation}
\begin{theorem}\label{thm:main}
The graph $H$ is connected and cubic, has odd girth exactly five, and
satisfies $\kappa(H)=4$. Consequently Hausmann's Question 5.4 has a
negative answer.
\end{theorem}

For comparison, Hausmann's Lemma 4.7, attributed there to Berke and
Szab\'o, gives a residual component bound of six in triangle-free
subcubic graphs. A different sufficient condition is the maximum-average-degree
bound of Cranston and Yancey \cite[Theorem 1, $k=3$]{cranston}:
$\operatorname{mad}(G)\leq18/7$ guarantees an independent set whose
complement is a forest with components of order at most three.
Since $H$ is cubic, $\operatorname{mad}(H)=3$; this hypothesis is absent
from Question 5.4. The counterexample concerns the proposed partition,
not the weighted maximum-cut inequality it was intended to imply.

\section{A complete classification proof}
Both layers in \eqref{eq:edges} are seven-cycles, and the spokes form
a matching. This proves simplicity, degree three, and connectedness.
A triangle within a layer is impossible. A triangle using both layers
would use two spokes sharing an endpoint, also impossible. The cycle
$a_0a_1a_2b_2b_0a_0$ has length five, so the odd girth is five.

\begin{lemma}\label{lem:classes}
Every maximal independent set of $H$ is equivalent to one of the five
sets in Table~\ref{tab:classes} under a simultaneous rotation or reflection
of the two index sets.
\end{lemma}
\begin{proof}
Write $A=\{i:a_i\in L\}$ and $B=\{i:b_i\in L\}$. Maximality means that
every vertex outside $L$ has a neighbor in $L$. Since the $a$-layer is a
seven-cycle, $|A|\leq3$.
If $A$ is empty, every $b_i$ must be selected to dominate $a_i$, which
contradicts independence. If $A=\{0\}$ after a rotation, domination of
$a_2,a_3,a_4,a_5$ forces $b_2,b_3,b_4,b_5\in L$, contradicting the edge
$b_2b_4$. Thus $|A|$ is two or three.

Suppose $|A|=2$. Rotation and reflection give $A=\{0,2\}$ or
$A=\{0,3\}$. In the first case, domination of $a_4,a_5$ forces
$b_4,b_5\in L$. Independence excludes $b_0,b_2,b_3,b_6$; domination
of $b_1$ then forces $b_1\in L$. This gives row III.
In the second case, $b_5$ is forced by $a_5$, while $b_0,b_3$ are
excluded. The other available inner vertices induce the path
$b_1b_6b_4b_2$. Their selected vertices must form a maximal independent
set of that path. For a path $x_1x_2x_3x_4$ the only such sets are
$\{x_1,x_3\}$, $\{x_1,x_4\}$, and $\{x_2,x_4\}$: select $x_1$ and
dominate the last edge, or omit $x_1$ and force $x_2,x_4$.
Thus $B$ is $\{1,4,5\}$, $\{1,2,5\}$, or $\{2,5,6\}$.
The reflection $i\mapsto3-i$ exchanges the first and third, giving
rows IV and V.

Suppose $|A|=3$. The cyclic gaps between its indices are at least two
and sum to seven, so they are $2,2,3$. Hence take $A=\{0,2,4\}$.
Every outer vertex is already dominated. The available inner vertices
induce the path $b_6b_1b_3b_5$, whose maximal independent sets are
$\{b_6,b_3\}$, $\{b_6,b_5\}$, and $\{b_1,b_5\}$.
The reflection $i\mapsto4-i$ exchanges the first and third.
These are rows I and II. Directly from \eqref{eq:edges}, each listed
set is independent and dominates every vertex outside it.
\end{proof}

\begin{table}[ht]
\centering
\small
\begin{tabular}{@{}clll@{}}
\toprule
Row & $\{i:a_i\in L\}$ & $\{i:b_i\in L\}$ & A path in $H-L$\\
\midrule
I   & $\{0,2,4\}$ & $\{1,5\}$   & $a_5a_6b_6b_4$\\
II  & $\{0,2,4\}$ & $\{5,6\}$   & $a_1b_1b_3a_3$\\
III & $\{0,2\}$   & $\{1,4,5\}$ & $a_3a_4a_5a_6$\\
IV  & $\{0,3\}$   & $\{1,2,5\}$ & $a_6a_5a_4b_4$\\
V   & $\{0,3\}$   & $\{1,4,5\}$ & $a_1a_2b_2b_0$\\
\bottomrule
\end{tabular}
\caption{The five maximal-independent-set classes and four connected
vertices outside each representative.}\label{tab:classes}
\end{table}

\begin{proof}[Proof of Theorem~\ref{thm:main}]
The graph properties were established above. Every independent set
$L$ extends to a maximal independent set $L^*$. Lemma~\ref{lem:classes}
and Table~\ref{tab:classes} supply a four-vertex path in $H-L^*$.
Rotations and reflections preserve \eqref{eq:edges}, so the same holds
for every representative's image. Since $H-L^*$ is an induced subgraph
of $H-L$, those four vertices remain in one component of $H-L$.
Thus $\kappa(H)\geq4$.

For the reverse inequality, take row II:
$L=\{a_0,a_2,a_4,b_5,b_6\}$. Its complementary components are precisely
\[
 \{a_1,b_1,b_3,a_3\},\qquad
 \{b_0,b_2,b_4\},\qquad
 \{a_5,a_6\}.
\]
Equation~\eqref{eq:edges} shows that these are paths of orders $4,3,2$,
with no edges between them. Hence $\kappa(H)\leq4$.
\end{proof}

\section{Consequences and the neighboring question}
\begin{corollary}\label{cor:universal}
Let $K_5$ be the least integer such that every finite simple subcubic
graph of odd girth five has an independent set whose complementary
components have order at most $K_5$. Then $4\leq K_5\leq6$.
\end{corollary}
\begin{proof}
Theorem~\ref{thm:main} gives the lower bound. Every graph in this class
is triangle-free, so the Berke--Szab\'o bound stated in
\cite[Lemma 4.7]{hausmann} gives the upper bound.
\end{proof}
Thus the existence of a constant, asked in the discussion of
\cite[Section 5, item 2]{aftd}, follows from that prior order-six bound.
The choices $K_5=4,5,6$ are not distinguished here.

\begin{corollary}\label{cor:induced}
If a finite graph $G$ contains $H$ as an induced subgraph, then
$\kappa(G)\geq4$. For every positive integer $t$ and every finite
bipartite graph $B$, $\kappa(tH\mathbin{\dot\cup}B)=4$.
\end{corollary}
\begin{proof}
For any independent $L$ of $G$, its restriction to $H$ is independent.
Four vertices connected in $H-L$ remain connected in $G-L$, proving
the first assertion. On a disjoint union, independent sets can be chosen
separately in each component, so $\kappa$ is the maximum of their
individual values. A bipartition class of $B$ leaves an independent
complement, of component order at most one. Apply Theorem~\ref{thm:main}
to each copy of $H$.
\end{proof}

Question 5.3 of \cite[Section 5]{hausmann} requires $L$ independent
and $G[V(C)\cup\{v\}]$ bipartite for every residual component $C$
and every $v\in L$. The printed statement has no separate requirement
that $G[R]$ be bipartite, whereas Theorem 2.2 of that paper does.
We allow empty parts in a vertex partition and distinguish the printed
quantifiers from their nonempty strengthening.

\begin{corollary}\label{cor:neighbor}
With empty parts allowed, the condition printed in Question 5.3 holds
for every graph by taking $L=\varnothing$. In contrast, $H$ admits no
such partition with $L\ne\varnothing$. In particular, $H$ admits no
such partition with the additional condition that $H[R]$ is bipartite.
\end{corollary}
\begin{proof}
The first assertion is vacuous universal quantification over $v\in L$.
Suppose that $L\ne\varnothing$ satisfies the condition. Then $H[R]$
is bipartite, since each residual component is an induced subgraph
of its union with any fixed $v\in L$.
Move any vertex $u\in R$ with no neighbor in $L$ into $L$.
Independence is preserved. Each new residual component $C'$ is contained
in an old one $C$. For each old $v\in L$, $H[V(C')\cup\{v\}]$ is
an induced subgraph of $H[V(C)\cup\{v\}]$; for the new vertex $u$,
it is an induced subgraph of the old bipartite graph $H[R]$.
Thus the condition persists until $L$ is maximal independent.

Use the five classes of Lemma~\ref{lem:classes}. In the table below,
each displayed odd cycle lies in $H[V(C)\cup\{v\}]$ for some
component $C$ of the corresponding complement and the indicated $v\in L$.
The repeated starting vertex is omitted. In row IV the cycle lies
entirely in $C$. Each edge follows from \eqref{eq:edges}.
\begin{center}
\small
\begin{tabular}{@{}cll@{}}
\toprule
Class & $v$ & Odd cycle\\
\midrule
I & $a_4$ & $a_4a_5a_6b_6b_4$\\
II & $a_2$ & $a_2a_1b_1b_3a_3$\\
III & $b_4$ & $b_4a_4a_5a_6b_6$\\
IV & $a_0$ & $a_4a_5a_6b_6b_4$\\
V & $a_0$ & $a_0a_1a_2b_2b_0$\\
\bottomrule
\end{tabular}
\end{center}
These five cycles contradict the condition in every maximal class.
For the last assertion, $L=\varnothing$ would leave the nonbipartite
graph $H$; a nonempty $L$ is excluded by the preceding argument.
\end{proof}

The nonempty and separately-bipartite variants are not asserted in the
printed Question 5.3. The corollary identifies their obstruction without
substituting them for the source's quantifiers.

\section{Coverage and verification}
\begin{center}
\small
\begin{tabular}{@{}>{\raggedright\arraybackslash}p{0.30\linewidth}>{\raggedright\arraybackslash}p{0.64\linewidth}@{}}
\toprule
Source item in \cite{hausmann} & Coverage\\
\midrule
Question 5.4 & Answered: no, by Theorem~\ref{thm:main}.\\
Earlier answer to Question 5.4 & Alternative proof; prior proof
AFTD (2026), using the Wagner graph.\\
Questions 5.1 and 5.2 & Not covered; these concern worst-case cut ratios.\\
Question 5.3 & Answered: yes with empty parts, by
Corollary~\ref{cor:neighbor}. Nonempty strengthening: not posed;
obstruction proved on $H$.\\
Question 5.5 & Not covered; $H$ has a five-cycle and is outside its
odd-girth-at-least-seven class.\\
Lemma 4.7 & Prior proof cited (Berke--Szab\'o, 2007, as stated by
Hausmann); order-six bound, not reproved here.\\
Conjecture 1.2 & Not covered; rejecting this structural sufficient
condition does not reject the weighted $4/5$ inequality.\\
Section 5, Wagner graph example & Prior proof cited (Hausmann, 2026);
an individual weighted cut ratio may decrease. Not reproved here.\\
Induced copies and disjoint unions of $H$ & Not posed; proved in
Corollary~\ref{cor:induced}.\\
\cite[Section 5, item 2]{aftd}, constant existence & Answered: yes;
prior order-six bound cited; Corollary~\ref{cor:universal} gives
$4\leq K_5\leq6$.\\
\cite[Section 5, item 3]{aftd}, order three at higher odd girth &
Not covered. The five-cycle in $H$ excludes this example.\\
\bottomrule
\end{tabular}
\end{center}

The written proof establishes the classification and the exact value
$\kappa(H)=4$. The accompanying \texttt{verify.py} enumerates all
$2^{14}=16{,}384$ vertex subsets: $435$ are independent, $49$ are
maximal independent, and none leaves all complementary components of
order at most three. It also checks Table~\ref{tab:classes} and the
partition attaining four.
It also checks every displayed odd cycle and tests the component-union
bipartiteness condition on all $435$ independent sets: the empty set
is its only solution on $H$.
An additional enumeration of $256$ subsets reproduces the earlier
Wagner-graph obstruction.

For the corresponding formal statement, the pinned
\href{https://github.com/the-omega-institute/trureturing/blob/6974dba01daa1a25cc74b87a86c191e95fe2c681/D5/S3/Combinatorics/Graph/ComponentPacking/HausmannQuestion54Refutation.lean}{\texttt{HausmannQuestion54Refutation.lean}}
defines \texttt{claim} as the full affirmative assertion of Question 5.4.
Its declaration \texttt{result :}\ $\neg$\texttt{claim} proves the negation using the
same labelled graph (labels $i=a_i$, $7+i=b_i$). Its helpers
\texttt{legal}, \texttt{four\_obstruction}, and \texttt{residual\_large}
establish, respectively, the degree/odd-girth hypotheses, four distinct
connected residual vertices for every independent set, and failure of
the component-order-three condition. The five-class classification and
the upper bound attaining four, and the three corollaries are written proofs
above. The odd-cycle table and the empty-set-only count are also finite
checks in \texttt{verify.py}; they are not additional Lean declarations.

The earlier Wagner-graph source is the
\href{https://github.com/rainwangphy/aftd-platform/blob/433d28c73eba86e5864ff92226afde4bb1c6e270/lean/AFTD/Kb/Tcs/SubcubicOddGirthFivePartitionComponentLeThreeNegation.lean}{pinned AFTD module}.
Its main declaration,
\nolinkurl{subcubic_odd_girth_five_partition_component_le_three_negation},
negates the universal partition assertion for triangle-free nonbipartite
subcubic graphs. The declaration ending in \texttt{\_wg\_key} gives
four connected complementary vertices for each independent set in the
Wagner graph. Its degree, triangle-exclusion, and nonbipartiteness
declarations end in \texttt{\_wg\_maxDegree}, \texttt{\_wg\_cliqueFree},
and \texttt{\_wg\_not\_bipartite}.
The explicit cycle $0,1,2,3,4,0$ in that graph supplies the exact
odd-girth-five hypothesis. These declarations concern the Wagner graph,
not the five-class proof or either new obstruction for $H$.

\begin{thebibliography}{9}
\bibitem{aftd}
AFTD, \emph{Subcubic Triangle-Free Non-Bipartite Graphs Do Not Always
Admit Independent Partitions with Complementary Component Order at Most Three},
November 5, 2026, especially Sections 3--5.
\href{https://rainwangphy.github.io/aftd-platform/blog/2026-11-05-op-1087/}{Wagner-graph counterexample}.

\bibitem{hausmann}
Will Hausmann, \emph{Weighted Max Cut in Subcubic Graphs},
arXiv:2611.05178v1, November 3, 2026.
\href{https://arxiv.org/html/2611.05178v1}{Full text},
especially Lemma 4.7 and Section 5, Questions 5.1--5.5.

\bibitem{cranston}
Daniel W. Cranston and Matthew P. Yancey,
\emph{Vertex Partitions into an Independent Set and a Forest with Each
Component Small}, SIAM Journal on Discrete Mathematics
\textbf{35}(3) (2021), 1769--1791.
\href{https://arxiv.org/html/2006.11445v2}{arXiv:2006.11445v2}, Theorem 1.
\end{thebibliography}

\smallskip
\noindent\small
\copyright\ 2026 Haobo Ma and Wenlin Zhang.
Licensed under \href{https://creativecommons.org/licenses/by/4.0/}
{Creative Commons Attribution 4.0 International (CC BY 4.0)}.
\end{document}
